% !TeX root=main.tex
% !TEX program = xelatex
% در این فایل، عنوان پایان‌نامه، مشخصات خود، متن تقدیمی‌، ستایش، سپاس‌گزاری و چکیده پایان‌نامه را به فارسی، وارد کنید.
% توجه داشته باشید که جدول حاوی مشخصات پروژه/پایان‌نامه/رساله و همچنین، مشخصات داخل آن، به طور خودکار، درج می‌شود.
%%%%%%%%%%%%%%%%%%%%%%%%%%%%%%%%%%%%
% دانشگاه خود را وارد کنید
\university{علم و صنعت ایران}
% دانشکده، آموزشکده و یا پژوهشکده  خود را وارد کنید
\faculty{دانشکده مهندسی کامپیوتر}
% گروه آموزشی خود را وارد کنید
% \department{گروه هوش مصنوعی و رباتیک}
% گروه آموزشی خود را وارد کنید
\subject{مهندسی کامپیوتر}
% عنوان پایان‌نامه را وارد کنید
\title{طراحی و پیاده‌سازی سامانه‌ی پرسش و پاسخ دوزبانه در حوزه‌ی اورولوژی مبتنی بر رویکرد تولید افزوده با بازیابی}

% نام استاد(ان) راهنما را وارد کنید
\firstsupervisor{دکتر ناصر مزینی}
%\secondadvisor{استاد مشاور دوم}

% مشخصات دانشجوی اول
\name{آیدا}
\surname{ریحانی}
\studentID{401...}

% مشخصات دانشجوی دوم
\secondname{صبا}
\secondsurname{باقری}
\secondstudentID{401521075}
% تاریخ پایان‌نامه را وارد کنید
\thesisdate{شهریور 1405}
% به صورت پیش‌فرض برای پایان‌نامه‌های کارشناسی تا دکترا به ترتیب از عبارات «پروژه»، «پایان‌نامه» و »رساله» استفاده می‌شود؛ اگر  نمی‌پسندید هر عنوانی را که مایلید در دستور زیر قرار داده و آنرا از حالت توضیح خارج کنید.
\projectLabel{پروژه}

% به صورت پیش‌فرض برای عناوین مقاطع تحصیلی کارشناسی تا دکترا به ترتیب از عبارات «کارشناسی»، «کارشناسی ارشد» و »دکترا» استفاده می‌شود؛ اگر  نمی‌پسندید هر عنوانی را که مایلید در دستور زیر قرار داده و آنرا از حالت توضیح خارج کنید.
\degree{کارشناسی}

\firstPage
\besmPage
\davaranPage

%\vspace{.5cm}
% در این قسمت اسامی اساتید راهنما، مشاور و داور باید به صورت دستی وارد شوند
%\renewcommand{\arraystretch}{1.2}
\begin{center}
\begin{tabular}{| p{8mm} | p{18mm} | p{.17\textwidth} |p{14mm}|p{.2\textwidth}|c|}
\hline
ردیف	& سمت & نام و نام خانوادگی & مرتبه \newline دانشگاهی &	دانشگاه یا مؤسسه &	امضـــــــــــــا\\
\hline
1 &     استاد راهنما				 & دکتر \newline ناصر مزینی & دانشیار & دانشگاه \newline علم و صنعت ایران & \\
\hline
\end{tabular}
\end{center}

\esalatPage
\mojavezPage


% چنانچه مایل به چاپ صفحات «تقدیم»، «نیایش» و «سپاس‌گزاری» در خروجی نیستید، خط‌های زیر را با گذاشتن ٪  در ابتدای آنها غیرفعال کنید.
 % پایان‌نامه خود را تقدیم کنید!

% \newpage
%\thispagestyle{empty}
%{\Large تقدیم به:}\\
%\begin{flushleft}
%{\huge
%همسر و فرزندانم\\
%\vspace{7mm}
%و\\
%\vspace{7mm}
%پدر و مادرم
%}
%\end{flushleft}


% سپاس‌گزاری
\begin{acknowledgementpage}
سپاس خداوندگار حکیم را که با لطف بی‌کران خود، آدمی را به زیور عقل آراست و ما را در مسیر کسب علم و دانش رهنمون ساخت.

در آغاز، وظیفهٔ خود می‌دانیم از زحمات بی‌دریغ، دلسوزانه و هدایت‌های علمی استاد راهنمای گران‌قدرمان، \textbf{جناب آقای دکتر ...} صمیمانه تشکر و قدردانی نماییم که بی‌تردید بدون راهنمایی‌های ارزنده و صبوری ایشان، به سرانجام رسیدن این پژوهش میسر نمی‌شد.

همچنین مراتب سپاس و قدردانی صمیمانهٔ خود را از \textbf{جناب آقای مهندس ...} که با دانش فنی، تجربیات گران‌بها و راهنمایی‌های ارزشمندشان ما را در پیشبرد و اجرای این پروژه صمیمانه یاری نمودند، ابراز می‌داریم.

در پایان، در برابر زحمات بی‌دریغ و مهر بی‌پایان پدران و مادران عزیزمان سر تعظیم فرود می‌آوریم و از خانواده‌های گرامی‌مان که همواره با دعای خیر، آرامش و پشتیبانی همه‌جانبه، انگیزه و گرمابخش مسیر ما در طول دوران تحصیل بودند، از صمیم قلب تشکر می‌کنیم.

% با استفاده از دستور زیر، امضای هر دو نفر درج می‌شود
\signature 
\end{acknowledgementpage}

%%%%%%%%%%%%%%%%%%%%%%%%%%%%%%%%%%%%
% کلمات کلیدی پایان‌نامه را وارد کنید
\keywords{تولید افزوده بازیابی \lr{(RAG)}، مدل‌های زبانی بزرگ \lr{(LLMs)}، اورولوژی، پردازش زبان طبیعی، \lr{LangGraph}}
%چکیده پایان‌نامه را وارد کنید، برای ایجاد پاراگراف جدید از \\ استفاده کنید. اگر خط خالی دشته باشید، خطا خواهید گرفت.
\fa-abstract{
پیشرفت مدل‌های زبانی بزرگ \lr{(LLMs)} چشم‌انداز جدیدی در پردازش زبان طبیعی و پاسخ‌گویی خودکار ایجاد کرده است. با این حال چالش توهم \lr{(Hallucination)} و محدود بودن پنجره پاسخ \lr{(Context window)} مانع استفاده از این مدل‌ها در حوزه‌های حساسی مانند پزشکی می‌شود. علاوه بر این، پشتیبانی کیفی از زبان فارسی در بسیاری از سیستم‌های تخصصی محدود است. برای غلبه بر این محدودیت‌ها و امکان بهره‌گیری مستقیم از کتب مرجع پزشکی به عنوان منبعی موثق، معماری تولید افزوده بازیابی \lr{(RAG)} به عنوان راهکار اصلی انتخاب شد. در این رساله، یک دستیار هوشمند اورولوژی دوزبانه (فارسی و انگلیسی) مبتنی بر این معماری توسعه یافته است.\\
در این راستا، متن کتب معتبر مرجع اورولوژی استخراج و پردازش شده و به عنوان پایگاه دانش در اختیار سیستم قرار گرفته است تا مدل به جای تکیه بر دانش درونی خود، پاسخ‌ها را از روی متون این کتب استخراج کند. سیستم پیشنهادی با استفاده از چارچوب \lr{LangGraph} به صورت چندعاملی \lr{(Multi-Agent)} طراحی شده و شامل عامل‌های مسیریابی \lr{(Router)}، ترجمه و بازیابی اطلاعات است. برای ذخیره و مدیریت محتوای کتب پزشکی از پایگاه داده برداری \lr{Qdrant} استفاده شده و مدل‌های زبانی منبع‌باز \lr{(HuggingFace)} در بستر \lr{Google Colab} مورد بهره‌برداری قرار گرفته‌اند. به منظور یکپارچه‌سازی، ارتباط بین هسته پردازشی و رابط کاربری \lr{(OpenWebUI)} از طریق یک \lr{API} سازگار با استاندارد \lr{OpenAI} در بستر \lr{FastAPI} پیاده‌سازی شده است.\\
نتایج نشان می‌دهد که این سیستم با جستجو در کتب مرجع، ارائه ارجاعات دقیق به منابع \lr{(Citations)} و نمایش مسیر تصمیم‌گیری \lr{(Breadcrumbs)}، علاوه بر رفع چالش‌های زبانی کاربر ایرانی، پاسخ‌هایی دقیق، شفاف و قابل اعتماد ارائه می‌دهد. در نهایت، کل سیستم با استفاده از \lr{Docker} کانتینرسازی شده و به صورت موفقیت‌آمیز مستقر \lr{(Deploy)} گردید تا استفاده عملیاتی و مقیاس‌پذیری آن در محیط‌های مختلف تضمین شود.
}


%\fa-abstract{
%%این پایان‌نامه، به بحث در مورد نوشتن پروژه، پایان‌نامه و رساله با استفاده از کلاس 
%%\lr{IUST-Thesis}
%%می‌پردازد. 
%%حروف‌چینی پروژه کارشناسی، پایان‌نامه یا رساله یکی از موارد پرکاربرد استفاده از زی‌پرشین است. 
%%زی‌پرشین بسته‌ای است که به همت آقای وفا خلیقی آماده شده است و امکان حروف‌چینی فارسی در  را  برای فارسی‌زبانان فراهم کرده است.
%%از جمله مزایای لاتک آن است که در صورت وجود یک کلاس آماده برای حروف‌چینی یک سند خاص مانند یک پایان‌نامه، کاربر بدون درگیری با جزییات حروف‌چینی و صفحه‌آرایی می‌توان سند خود را آماده نماید.
%
%شاید با قالب‌های لاتکی که برخی از مجلات برای مقالات خود عرضه می‌کنند مواجه شده باشید. اگر نظیر این کار در دانشگاههای مختلف برای اسناد متنوع آنها مانند پایا‌ن‌نامه‌ها آماده شود، دانشجویان به جای وقت گذاشتن روی صفحه‌آرایی مطالب خود، روی محتوای متن خود تمرکز خواهند نمود. به علاوه با آشنایی با لاتک خواهند توانست از امکانات بسیار این نرم‌افزار جهت نمایش بهتر دستآوردهای خود استفاده کنند.
%به همین خاطر، یک کلاس با نام 
%\lr{IUST-Thesis}
% برای حروف‌چینی پروژه‌ها، پایان‌نامه‌ها و رساله‌های دانشگاه علم و صنعت ایران با استفاده از نرم‌افزار زی‌پرشین،  آماده شده است. این فایل به 
%گونه‌ای طراحی شده است که کلیات خواسته‌های مورد نیاز  مدیریت تحصیلات تکمیلی دانشگاه علم و صنعت ایران را برآورده می‌کند و نیز، حروف‌چینی بسیاری از قسمت‌های آن، به طور خودکار انجام می‌شود.
%}

\abstractPage

\newpage\clearpage